\begin{textsample}{POS Dim 1 – to – Score 29.00 – to/to\_em\_5.txt}  \label{ex:f1_pos_280}
5.1 Língua Portuguesa O componente de Língua Portuguesa destina -se a oportunizar ao estudante o desempenho adequado da linguagem nas suas mais variadas situações e manifestações, até mesmo a estética, \textbf{uma} vez \textbf{que} o domínio da língua materna é fundamental ao acesso às demais áreas de conhecimento.
O ensino deve ser voltado para a função social da língua, requisito básico para o indivíduo construir seu processo de cidadania e, ainda, para participar autônoma e ativamente na sociedade.
Nessa perspectiva, o domínio linguístico \textbf{implica} interagir socialmente, dialogando em diferentes situações comunicativas, além de interpretar e produzir textos dos vários gêneros textuais \textbf{que} circulam em nossa sociedade.
Entende -se \textbf{que} o desenvolvimento da competência linguística integra : a leitura crítica dos diversos textos \textbf{que} circulam socialmente ; a produção escrita usando linguagem padrão ; a análise e uso da organização da estrutura da língua e a \textbf{percepção} das diferentes linguagens ( artística, corporal, verbal ).
Todas essas competências são elementos linguísticos, paralinguísticos e cinésicos como meio de compreensão do mundo e da produção, pois, \textbf{enquanto} se produz, é \textbf{preciso} examinar a própria produção, de modo \textbf{que} ela reflita pensamentos e sentimentos com a maior precisão possível.
Ensino Médio O ensino da Língua Portuguesa abrange a leitura, a escrita e a oralidade, além do discernimento de \textbf{que} todas as variantes de \textbf{uma} língua têm sua relevância, pois respondem às necessidades de cada povo e cada cultura.
É importante \textbf{que} o processo de interação entre os \textbf{sujeitos} possibilite a construção de sentidos e significados, considerando os diversos usos e funções da linguagem, bem como a \textbf{percepção} da relação entre a linguagem e a necessidade primeira de expressão de ideias, sentimentos e emoções vivenciadas pelo indivíduo.
Desse modo, concebe -se o domínio da Língua Portuguesa como fundamental ao desenvolvimento das competências de ouvir, falar, ler, escrever, discutir, \textbf{ver}, refletir e sentir.
Assim, a prática de ensino de \textbf{uma} língua deve partir da concepção \textbf{dialógica} de Bakhtin ( 1997 ), \textbf{que} aponta ser \textbf{inerente} a qualquer discurso o entrecruzamento deste com outros discursos.
Consequentemente, deve -se oportunizar ao estudante situações \textbf{que} favoreçam a sua \textbf{percepção} da coexistência \textbf{dialógica} dos diferentes gêneros discursivos, estabelecendo interações \textbf{que} vão além do texto, ou seja, \textbf{que} estão no interdiscurso.
Bakhtin ( 1986, p. 124 ) afirma \textbf{que} “ a língua \textbf{vive} e evolui \textbf{historicamente} na comunicação verbal concreta ”, não em sistema abstrato de formas linguísticas, e nem por \textbf{uma} enunciação monológica, mas se constitui por meio do “ fenômeno social da interação verbal realizada através da enunciação ou das enunciações.
A interação verbal constitui, assim, a realidade fundamental da língua ” ( BAKHTIN, 1986, p. 123 ).
Nessa perspectiva, cabe acrescentar as “ novas formas de ser, de se comportar, de discursar, de se relacionar, de se informar, de aprender ”, considerando \textbf{que} surgem “ novos tempos, novas tecnologias, novos textos, novas linguagens ” ( ROJO ; BARBOSA, 2015, p. 116 ), \textbf{fomentando} as atividades e ações diárias dos usuários da língua.
Na expectativa de linguagem como prática social e resultante da interação social, faz -se necessário entender e abordar “ a linguagem como \textbf{um} sistema sociossemiótico \textbf{que} permite escolhas e expressa significados em contextos linguísticos diversos ” ( HALLIDAY, 2004, p. 31 ).
Assim, estudar a linguagem é aprofundar o conhecimento de como as pessoas constroem significados por meio do uso da língua e de suas escolhas linguísticas nas várias situações, eventos, gêneros e atividades sociais em \textbf{que} participam, “ sejam estas mediadas por tecnologias ou produzidas em contextos profissionais, institucionais, em sala de aula, em \textbf{provas} de avaliação, entre outras ” ( REIS, 2010, p. 43 ).
É \textbf{preciso}, então, compreender e utilizar a linguagem em suas diferentes manifestações, como o lugar de constituição das relações sociais e, a partir daí, pensar o processo educacional.
O conhecimento deve ser compreendido como resultante de \textbf{um} processo histórico, construído socialmente no jogo das interações verbais.
Dessa forma, a sala de aula constitui -se em \textbf{um} espaço de interação, por meio de processos interlocutivos e variedades linguísticas para \textbf{que} os valores sociais a eles atribuídos possam ser discutidos e vivenciados.
Nesse sentido, Geraldi ( 1997, p. 135 ) propõe \textbf{que} “ o texto se constitua como ponto de partida e de chegada do ensino da língua ” e \textbf{que} se assuma a linguagem como atividade discursiva.
Desta forma, a leitura é entendida como \textbf{um} processo de interlocução, em \textbf{que} o leitor busca significações e reconstrói o texto, atribuindo -lhe sentido a partir de outras leituras, do seu conhecimento de mundo.
Segundo Paulo Freire ( 1996 ), é necessário valorizar o saber construído pelo \textbf{educando} na relação com a cultura local em \textbf{que} está inserido, pois é enaltecendo sua vivência e é na interação entre educador e \textbf{educando} \textbf{que} se estabelece o conhecimento.
No Ensino Médio, o Componente Curricular Língua Portuguesa deve ser oferecido nos três anos, em conformidade com a Lei nº 13.415/2017 ( BRASIL, 2017 ).
Diante disso, suas habilidades serão apresentadas \textbf{adiante}, organizadas por Campos de Atuação Social.
Por entender \textbf{que}, no Ensino Médio, os estudantes participam de forma significativa de diversas práticas sociais \textbf{que} envolvem a linguagem, pois, além de dominarem certos gêneros textuais/discursivos \textbf{que} circulam nos diferentes campos de atuação social considerados no Ensino Fundamental, eles já desenvolveram várias habilidades relativas aos usos das linguagens.
Conforme dispõe a BNCC : > Para orientar \textbf{uma} \textbf{abordagem} integrada dessas linguagens e de suas práticas, a área propõe \textbf{que} os estudantes possam vivenciar experiências significativas com práticas de linguagem em diferentes mídias ( impressa, digital, analógica ), situadas em campos de atuação social diversos, vinculados com o enriquecimento cultural próprio, as práticas cidadãs, o trabalho e a continuação dos estudos ( BRASIL, 2018, p. 485 ).
Assim, chegando ao Ensino Médio, os estudantes já têm a capacidade de participar significativamente de diversas práticas sociais mediadas pela linguagem.
Essa participação é possibilitada pelo fato de dominarem certos gêneros textuais/discursivos \textbf{que} circulam nas diferentes esferas sociais e de já terem desenvolvido várias habilidades relativas aos usos das linguagens.
Partindo dessa premissa, cabe ao componente curricular Língua Portuguesa aprofundar a análise sobre as linguagens e seus funcionamentos, intensificando a perspectiva analítica e crítica da leitura, escuta e produção de textos verbais e multissemióticos, e alargar as referências estéticas, éticas e políticas \textbf{que} cercam a produção e recepção de discursos, ampliando as possibilidades de fruição, de construção e produção de conhecimentos, de compreensão crítica e intervenção na realidade e de participação social dos jovens, nos âmbitos da cidadania, do trabalho e dos estudos ( BRASIL, 2018 ).
Desse modo, as práticas \textbf{contemporâneas} de linguagem ganham mais destaque no Ensino Médio, contemplando a cultura digital, as culturas juvenis, os novos letramentos e os multiletramentos, os processos colaborativos, as interações e atividades \textbf{que} têm lugar nas mídias e redes sociais, os processos de circulação de informações e a hibridização dos papéis nesse contexto ( de leitor/autor e produtor/consumidor ), exploradas no Ensino Fundamental.
Fenômenos como a pós-verdade e o efeito bolha, em função do impacto \textbf{que} produzem na fidedignidade do conteúdo disponibilizado nas redes, nas interações sociais e no trato com a diversidade, também são ressaltados ( BRASIL, 2018, p. 498 ). “ [... ] culturas de referência do \textbf{alunado} e de gêneros, mídias e linguagens por ele \textbf{conhecidos}, para buscar \textbf{um} enfoque crítico, pluralista, ético e democrático (... ) possibilitando a imersão em letramentos críticos ( ROJO, 2012, p. 21 ) ”.
Rojo ( 2012 ) explica ainda \textbf{que}, no multiletramento, converge -se a multiculturalidade, \textbf{que} é a variedade cultural das populações existentes nas sociedades, e a multimodalidade — variedade semiótica de composição dos textos pelos quais a multiculturalidade se comunica e informa.
Em relação à variedade cultural, Canclini ( 2008, p. 302 ) afirma \textbf{que} o \textbf{que} \textbf{vemos} à nossa volta são produções culturais letradas circulando socialmente, ou seja, “ \textbf{um} conjunto de textos híbridos de diferentes letramentos ( vernaculares e dominantes ), de diferentes campos ( \textbf{ditos} ‘ popular/de massa/erudito ’ ), caracterizados por escolha pessoal e política e de hibridização de produções de diferentes ‘ coleções ’ ”, constituindo o uso social da língua.
Nesse cenário de ensino e aprendizagem de Língua Portuguesa, a BNCC para o Ensino Médio indica o direito à aprendizagem e a progressão das habilidades, \textbf{que} acontecem ao se considerar : - a complexidade das práticas de linguagens e dos fenômenos sociais \textbf{que} repercutem nos usos da linguagem ( como a pós-verdade e o efeito bolha ) ; - a consolidação do domínio de gêneros do discurso/gêneros textuais já contemplados anteriormente e a ampliação do repertório de gêneros, sobretudo dos \textbf{que} \textbf{supõem} \textbf{um} grau maior de análise, síntese e reflexão ; - o aumento da complexidade dos textos lidos e produzidos em \textbf{termos} de temática, estruturação sintática, vocabulário, recursos estilísticos, orquestração de vozes e semioses ; - o foco maior nas habilidades envolvidas na reflexão sobre textos e práticas ( análise, avaliação, apreciação ética, estética e política, valoração, validação crítica, demonstração etc. ), considerando \textbf{que} as habilidades requeridas por processos de recuperação de informação ( identificação, reconhecimento, organização ) e por processos de compreensão ( comparação, distinção, estabelecimento de relações e inferência ) já foram desenvolvidas no Ensino Fundamental ; - a atenção maior nas habilidades envolvidas na produção de textos multissemióticos mais analíticos, críticos, propositivos e criativos, abarcando sínteses mais complexas, produzidos em contextos \textbf{que} \textbf{suponham} apuração de fatos, curadoria, levantamentos e pesquisas e \textbf{que} possam ser vinculados de forma significativa aos contextos de estudo/construção de conhecimentos em diferentes áreas, a experiências estéticas e produções da cultura digital e à discussão e proposição de ações e projetos de relevância pessoal e para a comunidade ; - o incremento da consideração das práticas da cultura digital e das culturas juvenis globais, regionais e locais, por meio do aprofundamento da análise de suas práticas e produções culturais em circulação, de \textbf{uma} maior incorporação de critérios técnicos e estéticos na análise e autoria das produções e vivências mais intensas de processos de produção colaborativos ; - a ampliação de repertório, considerando a diversidade cultural, de maneira a abranger produções e formas de expressão diversas – literatura juvenil, literatura periférico-marginal, o culto, o clássico, o \textbf{popular}, cultura de massa, cultura das mídias, culturas juvenis etc.
E em suas múltiplas repercussões e possibilidades de apreciação, em processos \textbf{que} envolvem adaptações, remidiações, estilizações, paródias, HQs, minisséries, filmes, videominutos, games etc. ; - a inclusão de obras da tradição \textbf{literária} brasileira e de suas referências ocidentais – em especial da literatura portuguesa, assim como obras mais complexas da literatura \textbf{contemporânea}, africana e latino-americana ( BRASIL, 2018, p. 499-500 ).
Além da indicação da progressão das habilidades, ressalta -se \textbf{aqui} a importância de \textbf{que} o estudo e a inclusão da literatura tocantinense ocorram de forma gradativa, iniciando pelos \textbf{autores} em cada localidade, abrangendo posteriormente os regionais e os estaduais.
O objetivo é promover, ressaltar e valorizar a literatura local e regional de forma ampla, potencializando o conhecimento cultural, artístico e sociolinguístico adquiridos pelas leituras das obras \textbf{literárias} tocantinenses.

% matched lemmas: abordagem, adiante, alunado, aqui, autor, conhecido, contemporâneo, dialógico, dizer, educar, enquanto, fomentar, historicamente, implicar, inerente, literário, percepção, popular, preciso, prova, que, sujeito, supor, termos, um, ver, viver
\end{textsample}
